\begingroup
\renewcommand{\endProof}{\par}
\begin{Proof}
Let $h\in(0,R(1-\lambda)/2)$ and fix an index $k=1,\ldots,n-1$ (i.e., $k\ne n$). Now choose $\zeta\in C_b^\infty(D^+)$ such that\begin{enumerate}\item $0\leq\zeta\leq1$,\item $\zeta\equiv1$ on $Q^+$,\item $U:=\operatorname{supp}\zeta\subset B_{R(1+\lambda)/2}(0)$.\end{enumerate}\textbf{Note:} The function $\zeta$ is not in $\mathcal{D}(D^+)$ since $\zeta\not\equiv0$ on the flat part of the boundary of $D^+$: $\{x_n=0\}$.

Now define
\begin{equation}\label{eq:flat-boundary-tangential-difference-test-function}
v:=-D_k^{-h}(\zeta^2D_k^hu).
\end{equation}
After some manipulations using the definition of the difference quotients, we get the following identity.
\begin{equation}\label{eq:flat-boundary-test-function-translation-expansion}
\begin{aligned}
v(\bm{x})=-\frac{1}{h^2}\Bigl(&\zeta^2(\bm{x})[u(\bm{x}+h\bm{e}_k)-u(\bm{x})]\\
&+\zeta^2(\bm{x}-h\bm{e}_k)[u(\bm{x}-h\bm{e}_k)-u(\bm{x})]\Bigr).
\end{aligned}
\end{equation}
Note that in constructing $v$ we have used translations only in directions tangential to the plane $x_n=0$. The key idea is that we can ``slide the support of $u$\textquotesingle\textquotesingle\ along the plane $x_n=0$ without destroying the boundary conditions. Also note that none of the translations moves the support of $\zeta$ outside of $\overline{D^+}$. These facts ensure that\begin{enumerate}\item $v$ is well defined on $D^+$,\item $v\in H^1(D^+)$ (since $u\in H^1(D^+)$),\item $v\in H_0^1(D^+)$ (since $\zeta$ is zero on the curved part of the boundary of $D^+$ and $u$ is zero on the flat part of the boundary (and the same goes for any of the translations of $\zeta$ and $u$)).\end{enumerate}We define
\begin{equation}\label{eq:flat-boundary-localized-tangential-difference}
w:=\zeta D_k^hu,
\end{equation}
so that $v=-D_k^{-h}(\zeta w)$.

Now, we follow a procedure similar to the derivation of the estimate \eqref{eq:interior-regularity-shifted-principal-form-bound} of Theorem~\ref{thm:second-order-interior-regularity}, but using $v$ as defined above. We get\[\begin{aligned}
-\int_{D^+}gD_k^{-h}(\zeta w)\,\mathrm{d}\bm{x}
&=\int_{D^+}a_{ij}u_{x_j}v_{x_i}\,\mathrm{d}\bm{x}\\
&=\int_{D^+}D_k^h(a_{ij}u_{x_j})(\zeta w)_{x_i}\,\mathrm{d}\bm{x}\\
&=\int_{D^+}a_{ij}(\bm{x}+h\bm{e}_k)D_k^hu_{x_j}(\zeta w)_{x_i}\,\mathrm{d}\bm{x}\\
&\quad+\int_{D^+}D_k^ha_{ij}u_{x_j}(\zeta w)_{x_i}\,\mathrm{d}\bm{x}\\
&=\int_{D^+}a_{ij}(\bm{x}+h\bm{e}_k)w_{x_j}w_{x_i}\,\mathrm{d}\bm{x}\\
&\quad-\int_{D^+}a_{ij}(\bm{x}+h\bm{e}_k)\zeta_{x_j}D_k^huw_{x_i}\,\mathrm{d}\bm{x}\\
&\quad+\int_{D^+}a_{ij}(\bm{x}+h\bm{e}_k)D_k^hu_{x_j}\zeta_{x_i}w\,\mathrm{d}\bm{x}\\
&\quad+\int_{D^+}D_k^ha_{ij}u_{x_j}(\zeta w)_{x_i}\,\mathrm{d}\bm{x}.
\end{aligned}\]Here $g$ is defined as in \eqref{eq:interior-regularity-reduced-right-hand-side}. Rearranging, we get
\begin{equation}\label{eq:flat-boundary-principal-energy-four-term-decomposition}
\int_{D^+}a_{ij}(\bm{x}+h\bm{e}_k)w_{x_j}w_{x_i}\,\mathrm{d}\bm{x}=I_1+I_2+I_3+I_4,
\end{equation}
where
\begin{equation}\label{eq:flat-boundary-energy-term-one}
I_1:=-\int_{D^+}gD_k^{-h}(\zeta w)\,\mathrm{d}\bm{x},
\end{equation}

\begin{equation}\label{eq:flat-boundary-energy-term-two}
I_2:=\int_{D^+}a_{ij}(\bm{x}+h\bm{e}_k)\zeta_{x_j}D_k^huw_{x_i}\,\mathrm{d}\bm{x},
\end{equation}

\begin{equation}\label{eq:flat-boundary-energy-term-three}
I_3:=-\int_{D^+}a_{ij}(\bm{x}+h\bm{e}_k)D_k^hu_{x_j}\zeta_{x_i}w\,\mathrm{d}\bm{x},
\end{equation}

\begin{equation}\label{eq:flat-boundary-energy-term-four}
I_4:=-\int_{D^+}D_k^ha_{ij}u_{x_j}(\zeta w)_{x_i}\,\mathrm{d}\bm{x}.
\end{equation}
We estimate these terms using techniques which should be familiar to the reader from the proof of Theorem~\ref{thm:second-order-interior-regularity} (basically Hölder\textquotesingle s inequality and Lemma~\ref{lem:difference-quotient-weak-derivative-bound}). We will make use of the estimate
\begin{equation}\label{eq:flat-boundary-test-function-l2-bound}
\|v\|_2=\|D_k^{-h}(\zeta w)\|_2\leq2\|\zeta\|_{1,\infty}\|w\|_{1,2}.
\end{equation}

\end{Proof}
\endgroup
